\documentclass{article}

%!TEX root = ../main.tex

\usepackage{url}
\def\UrlFont{\rmfamily}

\usepackage{enumitem}
% [noitemsep,topsep=0pt,labelwidth=0pt,listparindent=-30pt,leftmargin=30pt,rightmargin=30pt]
\usepackage{cite}
\usepackage{amsmath,amssymb,amsfonts}
\usepackage{algorithmic}
\usepackage{graphicx}
\graphicspath{{./img/}}

\usepackage{textcomp}
\usepackage{xcolor}
\usepackage{booktabs}
\usepackage{caption}
\usepackage{xspace}
\usepackage{float}
\usepackage{marginnote}
\usepackage{tabularx}
\usepackage{changepage, threeparttable}
\usepackage{attrib}
\usepackage{svg}
\usepackage{multirow}


\graphicspath{{./img/}}
%!TEX root = ../main.tex


\newcommand\cc[1]{\textcolor{red}{#1}}
\newcommand\tobedeleted[1]{\textcolor{red}{#1}}



\title{Individual hand-in for the exam session %
	\footnote{%
        \tobedeleted{\newline %
        The submitted document must keep the structure presented in this template. \\ %
        Limitation on the length of content are stated in red. \\ %
	The supporting text, which is marked in red, must be removed. \\ %
	Finally, remember to submit only the pdf and to change the file name to: \\ %
	``\texttt{<Group Number>\char`_<Student Username>\char`_<Hand-in Number>.pdf}'', e.g., %
	\texttt{XX\char`_pate\char`_3.pdf}}} %
	\\[10pt] %
	\normalsize{[v26.1]}}
 
\date{\today}
\begin{document}
\maketitle

\nocite{2001:manifesto}

\vfill
\begin{description}[noitemsep,topsep=0pt,labelwidth=0pt,listparindent=-30pt,leftmargin=30pt,rightmargin=30pt]
\item[Team name: ] Group 16 
\item[Company: ] Alfa Laval  
\item[Teaching assistant: ] Hiep Cao
\item[Agile coach: ] Henning Therkelsen
\end{description}
\begin{table}[h]\begin{center}
\begin{tabular*}{1\textwidth}{@{}ccp{8cm}}
\toprule
Role & email & Name \\
\midrule
PO & Kasper.Haarslov@alfalaval.com & Kasper Haarsløv \\
Specialist & morten.gronkjaer@alfalaval.com & Morten Groenkjaer \\
SM & maev@itu.dk & Maria Melina Evdaimon \\
dev & cmje@itu.dk & Casper Meier Jenkins \\
dev & ebll@itu.dk & Emilie Bliddal Ravn Larsen \\
dev & jpre@itu.dk & Joakim-David Stauning \\
dev & kfje@itu.dk & Kasper Fjord Grønvold Jensen \\
dev & mhil@itu.dk & Milja Noomi Hildigunnsdottir \\
dev & moto@itu.dk & Morgan Møller Torp \\
dev & merl@itu.dk & Morten Holtorp Erlandsen \\
dev & thsch@itu.dk & Therese Schultz-Nielsen \\
dev & vmol@itu.dk & Vincent Jørss Molenaar\\
\bottomrule
\end{tabular*}\end{center}\end{table}


\newpage


\section{Tasks}
\begin{itemize}
    \item Execute the test at \url{https://www.scrum.org/open-assessments/scrum-open}.
    \item Percentage: 76.7\% 
    \item "When does a Developer become the sole owner of an item on the Sprint Backlog?" - I just thought of this question differently than what was intended. In my mind it makes sense to delegate taskes out to different people so the developers know who is working on it.
\end{itemize}


\section{Reflect on one main point you learned in the last two sprints. This must be stated as an educational observation related to Scrum or other work practices (e.g., what did you learn, how). Include at least one concrete reference to an artifact from Sprints 5–6 (e.g., backlog item, PR/commit, sprint event, retro note, TA/coach feedback).}
\label{sec:actual}

I learned that it can be hard to communicate over text, and that ones intentions might not be so clearly defined over text. I had an incident were I reviewed a pr and the author of the pr read the reviewed as me yelling at them. This caused a bit of drama and was on me as I should have communicated better.

I also learned that it might be a good idea for when multiple people review a pr together that they then add a little comment telling the author who reviewed it. This was because I had made a few small changes that I felt was ignored, since it looked like nobody reviewed it. I latter learned that multiple people had reviewed it together but was not transparent about it. When I learned that I felt much better since I learned that I was not ignored



\section{During the experience exchange or community of practice sessions (MSc), (1) how has my active participation benefited others and (2) what did I discuss with others that helped me?}
\label{sec:exchange}
The biggest thing that I learned was that everybody else felt like they could not deliver the finished product that the PO's wanted. This was nice since there were a lot less stress about not finishing a complete product. 
A complete product does not mean that we have nothing to show. It just mean that the ideal dream product could not be finished

\section{Final reflections on the project and learning}

\tobedeleted{This is the most important part of the final hand-in.}

\tobedeleted{Your reflection must:}

\begin{itemize}
    \item \tobedeleted{for BSc students: be between 1 and 2 pages,}
    \item \tobedeleted{for MSc students: be between 2 and 3 pages and include references to relevant 
  research, as required by the ILOs.}
\end{itemize}

\tobedeleted{Your reflection must reference at least three concrete project artifacts, such as:}
\begin{itemize}
    \item \tobedeleted{backlog items (include IDs),}
    \item \tobedeleted{PRs or commits (include identifiers),}
    \item \tobedeleted{retro board items,}
    \item \tobedeleted{sprint goals,}
    \item \tobedeleted{teacher, TA or agile coach feedback,}
    \item \tobedeleted{design or architectural changes.}   
\end{itemize}


\tobedeleted{Structure your reflection around the following:}

\begin{itemize}
    \item \tobedeleted{Early sprints (1–2): what surprised you and why.}
    \item \tobedeleted{Middle sprints (3–4): one challenge or failure, why it happened, and what you 
   learned.}
    \item \tobedeleted{Late sprints (5–6): what improved, what changed, and why.}
    \item \tobedeleted{Collaboration and teamwork: describe one conflict, disagreement, or tension 
   (use initials, not names), how it was addressed, and what you learned.}
    \item \tobedeleted{Scrum/SE practices: what you would do differently next time and why.}
    \item \tobedeleted{MSc only: relate your learning to one or two relevant research papers or 
   frameworks, and explain how they helped you understand or analyse your 
   project experience.}   
\end{itemize}


\tobedeleted{Generic or hypothetical descriptions are insufficient. Your reflection must be 
tied to your actual experience and supported by concrete evidence.}














\newpage
\bibliographystyle{plain}
\bibliography{misc/cite}

\end{document}
